\Article{363}
The main task is to determine these remainders. First, find the remainder that the largest period gives on division by the product of all the appropriate prime functions. If this product is
\[
\equiv x^\lambda-Ax^{\lambda-1}+\text{etc.}
\]
then this remainder is $\equiv A$. The form of this product is readily found, and Chapter \textsc{vi} shows that $A=0$ if $\nu$ is divisible by a square; otherwise $A$ is either $=+1$ or $=-1$, according as the number of prime factors of $\nu$ is even or odd.

Now split this largest period into smaller periods, and represent the terms of any one period by $x^{kp^\pi u}$, so that $k$ is a number
\SourceLeaf{242}
fixed within each period but varying between periods, while $\pi$ and $u$ vary within each period and can take in the other periods the same values that they take in any one period. For the moment, take some prime function $P$ as the base, and let the periods $\Sigma x^{p^\pi u}$, $\Sigma x^{k'p^\pi u}$, etc., on division by it, give the remainders $A,A'$, etc., respectively. Then $\Sigma x^{p^\pi u}-A$ is divisible by the product of all the functions $(P,\rho^u)$; $\Sigma x^{k'p^\pi u}-A'$ by the product of all the functions $(P,\rho^{k'u})$; etc., etc. It is readily seen that the quantities $A,A'$, etc., are roots of a known congruence. Namely, let the periods of the roots of the equation $x^\nu=1$ corresponding to the preceding periods be roots of the equation $Q=0$; then $A,A'$, etc., will be roots of the congruence $Q\equiv0$. For we have
\[
\begin{aligned}
A+A'+\text{etc.}&\equiv\text{the sum of the periods},\\
AA+A'A'+\text{etc.}&\equiv\text{the sum of the squares of the periods}
\end{aligned}
\]
etc., etc. The calculation is entirely similar to the one we set out in Chapter \textsc{vi}, if in place of $\rho$ we substitute $x$, since here too we can put unity for $x^\nu$, as we did there for $\rho^\nu$.

Once the roots $A,A'$, etc., have been found, choose one as the remainder of the period $\Sigma x^{p^\pi u}$, and use it to order the remainders of the others in the same way as in Chapter \textsc{vi}. This choice is left free here too, since the function $P$ is still entirely undetermined. The subsequent calculation is wholly analogous to the one we treated in Chapter \textsc{vi}; setting out every detail would take us too far. Finally, once we have reached $\Sigma x^{p^\pi}$, the main task is complete. For, on putting
\[
P\equiv x^m+ax^{m-1}+bx^{m-2}+\text{etc.}
\]
we have $-a\equiv\Sigma x^{p^\pi}$. The second coefficient of the other functions $(P,\rho^k)$ is obtained in the same way, and from this the remaining coefficients of $P$ can be determined. It may often happen that we reach identical congruences from which nothing seems to follow. How this difficulty can be met will be shown below.

\Article{364}
An example will make all this much clearer. The proposed task is to resolve the function $x^{15}-1$ into factors modulo $17$. We have $m=4$, and since the product of all the elementary functions is
\[
\equiv\frac{x^{15}-1\pdot x-1}{x^3-1\pdot x^5-1}=x^8-x^7+x^5-x^4+x^3-x+1
\]

\SourceLeaf{243}
there will therefore be only two prime factors of degree four, $P$ and $P'$. Now distribute $x,xx,x^4,x^7,x^8,x^{11},x^{13},x^{14}$ into these two periods:
\[
\begin{aligned}
\Sigma x^{17^\alpha}&\equiv x+xx+x^4+x^8,\\
\Sigma x^{7\pdot17^\alpha}&\equiv x^7+x^{11}+x^{13}+x^{14}
\end{aligned}
\]
Modulo one of the functions $P,P'$, let
\[
\Sigma x^{17^\alpha}\equiv A,\quad\Sigma x^{7\pdot17^\alpha}\equiv A'
\]
Then
\[
\begin{aligned}
A+A'&\equiv1\\
AA&\equiv\Sigma x^{2\pdot17^\alpha}+\Sigma x^{3\pdot17^\alpha}+\Sigma x^{5\pdot17^\alpha}+\Sigma x^{9\pdot17^\alpha}\\
A'A'&\equiv\Sigma x^{14\pdot17^\alpha}+\Sigma x^{6\pdot17^\alpha}+\Sigma x^{5\pdot17^\alpha}+\Sigma x^{3\pdot17^\alpha}
\end{aligned}
\]
and hence
\[
\begin{aligned}
AA+A'A'&\equiv\Sigma x^{17^\alpha}+\Sigma x^{7\pdot17^\alpha}+4\Sigma x^{3\pdot17^\alpha}+2\Sigma x^{5\pdot17^\alpha}\\
&\equiv1-4-4\equiv-7
\end{aligned}
\]
Thus $A$ and $A'$ will be the roots of the congruence
\[
xx-x+4\equiv0\;(\mathrm{mod.}\,17)
\]
which are $6,12$. It follows that $P$ divides
\[
x^8+x^4+xx+x-6
\]
and is
\[
\equiv x^4-6x^3-2xx-12x+1
\]
Further, $P'$ will be $\equiv(P,\rho^7)$, and will be
\[
\equiv x^4-12x^3-2xx-6x+1
\]
